\chapter{عبارات منظم}

در حساب با عملگرهای $+$ و $\times$ عبارت‌هایی مثل $(5+3)\times4$ می‌سازیم که مقدارشان یک عدد است. به همین شکل، با عملیات منظم عبارت‌هایی می‌سازیم که مقدارشان یک \textbf{زبان} است. به این عبارت‌ها \term{عبارات منظم}{Regular Expressions} می‌گوییم. عبارات منظم روشی فشرده و قدرتمند برای توصیف الگوها هستند و در برنامه‌نویسی (مثلاً جست‌وجوی متن و ساخت کامپایلر) کاربرد زیادی دارند.

\section{تعریف و مثال}

مثلاً عبارت $(0\cup1)0^*$ کوتاه‌شده $(\{0\}\cup\{1\})\circ\{0\}^*$ است: مجموعه همه رشته‌هایی که با یک $0$ یا یک $1$ شروع می‌شوند و بعد هر تعداد $0$ دارند. چند قرارداد:
\begin{itemize}
\item نماد الحاق $\circ$ معمولاً نوشته نمی‌شود.
\item \textbf{اولویت عملگرها:} اول ستاره، بعد الحاق، بعد اجتماع (مثل توان، ضرب و جمع در حساب). پس $01^*\cup1$ یعنی $\big(0(1^*)\big)\cup1$.
\item اگر $\Sigma=\{0,1\}$، خود $\Sigma$ را به عنوان عبارت $(0\cup1)$ به کار می‌بریم. پس $\Sigma$ همه رشته‌های به طول یک، $\Sigma^*$ همه رشته‌ها، $0\Sigma^*$ رشته‌هایی که با $0$ شروع می‌شوند و $\Sigma^*1$ رشته‌هایی که به $1$ ختم می‌شوند را توصیف می‌کند.
\end{itemize}

\begin{definition}[عبارت منظم]
$R$ یک \textbf{عبارت منظم} است اگر یکی از این‌ها باشد:
\begin{enumerate}
\item $a$ برای یک $a\in\Sigma$؛ زبان آن $\{a\}$ است.
\item $\eps$؛ زبان آن $\{\eps\}$ است که فقط یک رشته (رشته تهی) دارد.
\item $\emptyset$؛ زبان آن تهی است و هیچ رشته‌ای ندارد.
\item $(R_1\cup R_2)$، $(R_1\circ R_2)$ یا $(R_1^*)$ که $R_1$ و $R_2$ عبارت منظم‌اند.
\end{enumerate}
زبان توصیف‌شده با $R$ را با $L(R)$ نشان می‌دهیم. همچنین $R^+=RR^*$ و $R^k$ الحاق $k$ نسخه از $R$ است.
\end{definition}

\begin{note}[تعریف دوری یا استقرایی؟]
ممکن است به نظر برسد عبارت منظم را با خودش تعریف کرده‌ایم (تعریف دوری). اما این‌طور نیست: $R_1$ و $R_2$ همیشه \textbf{کوچک‌تر} از $R$ هستند. پس این یک \textbf{تعریف استقرایی} است که از عبارت‌های پایه (موارد ۱ تا ۳) شروع می‌کند و عبارت‌های بزرگ‌تر را از کوچک‌ترها می‌سازد.
\end{note}

\begin{example}
روی $\Sigma=\{0,1\}$:
\begin{ltrfig}
\renewcommand{\arraystretch}{1.25}
\begin{tabular}{ll}
$0^*10^*$ & $\{w\st w$ contains a single $1\}$\\
$\Sigma^*1\Sigma^*$ & $\{w\st w$ has at least one $1\}$\\
$\Sigma^*001\Sigma^*$ & $\{w\st w$ contains $001$ as a substring$\}$\\
$1^*(01^+)^*$ & $\{w\st$ every $0$ in $w$ is followed by at least one $1\}$\\
$(\Sigma\Sigma)^*$ & $\{w\st |w|$ is even$\}$\\
$(\Sigma\Sigma\Sigma)^*$ & $\{w\st |w|$ is a multiple of $3\}$\\
$01\cup10$ & $\{01,10\}$\\
$0\Sigma^*0\cup1\Sigma^*1\cup0\cup1$ & $\{w\st w$ starts and ends with the same symbol$\}$\\
$(0\cup\eps)1^*$ & $01^*\cup1^*$\\
$(0\cup\eps)(1\cup\eps)$ & $\{\eps,0,1,01\}$\\
$1^*\emptyset$ & $\emptyset$\\
$\emptyset^*$ & $\{\eps\}$
\end{tabular}
\end{ltrfig}
\end{example}

دو مورد آخر ارزش توضیح دارند. الحاق هر زبانی با $\emptyset$ تهی است، چون برای ساختن رشته $xy$ باید $y$ را از زبان تهی برداریم که ممکن نیست. اما $\emptyset^*=\{\eps\}$، چون ستاره همیشه الحاق \emph{صفر} رشته (یعنی $\eps$) را دارد.

\subsection{اتحادهای ساده}
برای هر عبارت منظم $R$:
\[R\cup\emptyset=R,\qquad R\circ\eps=R.\]
اما دو عبارت شبیه به این‌ها در حالت کلی درست \textbf{نیستند}:
\begin{itemize}
\item $R\cup\eps\neq R$: اگر $R=0$ آنگاه $L(R)=\{0\}$ ولی $L(R\cup\eps)=\{0,\eps\}$.
\item $R\circ\emptyset\neq R$: اگر $R=0$ آنگاه $L(R)=\{0\}$ ولی $L(R\circ\emptyset)=\emptyset$.
\end{itemize}

\begin{example}[کاربرد در کامپایلرها]
بخش «تحلیل لغوی» کامپایلر کد برنامه را به \textbf{توکن}ها (مثل نام متغیر و عدد) می‌شکند و معمولاً از روی عبارات منظم به‌طور خودکار ساخته می‌شود. مثلاً اگر $D=\{0,1,\dots,9\}$ مجموعه ارقام باشد، عددهای اعشاری علامت‌دار با این عبارت توصیف می‌شوند:
\[(+\cup-\cup\eps)\,(D^+\cup D^+.D^*\cup D^*.D^+).\]
رشته‌های $72$، $3.14159$، $+7.$ و $-.01$ همه با این عبارت تطبیق دارند. (به اولویت عملگرها دقت کنید.)
\end{example}

\section{هم‌ارزی عبارات منظم و ماشین‌های متناهی}

\begin{theorem}\label{thm:re-equiv}
یک زبان منظم است اگر و تنها اگر یک عبارت منظم آن را توصیف کند.
\end{theorem}

این قضیه دو جهت دارد که هر کدام را در یک لم ثابت می‌کنیم.

\subsection{از عبارت منظم به \lr{NFA}}

\begin{lemma}\label{lem:re-nfa}
اگر یک عبارت منظم زبانی را توصیف کند، آن زبان منظم است.
\end{lemma}

\begin{proofidea}
عبارت منظم $R$ را به یک \lr{NFA} تبدیل می‌کنیم. چون \lr{NFA} و \lr{DFA} هم‌ارزند، این کافی است. تبدیل را روی ساختار استقرایی تعریف عبارت منظم انجام می‌دهیم: برای عبارت‌های پایه ماشین می‌سازیم و برای عبارت‌های مرکب از ساخت‌های بستاری فصل ۵ استفاده می‌کنیم.
\end{proofidea}

\begin{proof}
شش حالت تعریف عبارت منظم را بررسی می‌کنیم:
\begin{enumerate}
\item $R=a$ برای $a\in\Sigma$، پس $L(R)=\{a\}$. ماشین $N=(\{q_1,q_2\},\Sigma,\delta,q_1,\{q_2\})$ با $\delta(q_1,a)=\{q_2\}$ و $\delta(r,b)=\emptyset$ برای بقیه حالت‌ها.
\item $R=\eps$، پس $L(R)=\{\eps\}$. ماشین $N=(\{q_1\},\Sigma,\delta,q_1,\{q_1\})$ با $\delta(q_1,b)=\emptyset$ برای هر $b$.
\item $R=\emptyset$. ماشین $N=(\{q_1\},\Sigma,\delta,q_1,\emptyset)$ با $\delta(q_1,b)=\emptyset$ برای هر $b$.
\end{enumerate}
\begin{automaton}[node distance=1.6cm]
  \node[state,initial] (a1) {$q_1$};
  \node[state,accepting,right=of a1] (a2) {$q_2$};
  \path (a1) edge node {$a$} (a2);
  \node[state,initial,accepting,right=2.4cm of a2] (b1) {$q_1$};
  \node[state,initial,right=2.6cm of b1] (c1) {$q_1$};
  \node[below=4mm of $(a1)!0.5!(a2)$,font=\small] {case 1};
  \node[below=4mm of b1,font=\small] {case 2};
  \node[below=4mm of c1,font=\small] {case 3};
\end{automaton}
\begin{enumerate}[resume]
\item $R=R_1\cup R_2$،
\item $R=R_1\circ R_2$،
\item $R=R_1^*$.
\end{enumerate}
در سه حالت آخر، طبق فرض استقرا برای $R_1$ و $R_2$ (که کوچک‌ترند) \lr{NFA} داریم و با ساخت‌های اثبات بستاری تحت اجتماع، الحاق و ستاره (فصل ۵) یک \lr{NFA} برای $R$ می‌سازیم.
\end{proof}

\begin{example}[تبدیل $(ab\cup a)^*$ به \lr{NFA}]
ماشین را از درون به بیرون می‌سازیم:
\begin{automaton}[node distance=1.15cm,every state/.style={draw=jzMain,fill=jzMain!7,minimum size=6mm,inner sep=0pt}]
  \node (la) {$a$:};
  \node[state,right=0.3cm of la] (a1) {}; \node[state,accepting,right=of a1] (a2) {};
  \path (a1) edge node {$a$} (a2);
  \node[right=1.4cm of a2] (lb) {$b$:};
  \node[state,right=0.3cm of lb] (b1) {}; \node[state,accepting,right=of b1] (b2) {};
  \path (b1) edge node {$b$} (b2);
  \node[below=1.1cm of la] (lab) {$ab$:};
  \node[state,right=0.3cm of lab] (c1) {}; \node[state,right=of c1] (c2) {}; \node[state,right=of c2] (c3) {}; \node[state,accepting,right=of c3] (c4) {};
  \path (c1) edge node {$a$} (c2) (c2) edge node {$\eps$} (c3) (c3) edge node {$b$} (c4);
\end{automaton}
\begin{automaton}[node distance=1.15cm,every state/.style={draw=jzMain,fill=jzMain!7,minimum size=6mm,inner sep=0pt}]
  \node (l) {$ab\cup a$:};
  \node[state,right=0.3cm of l] (s) {};
  \node[state,above right=0.5cm and 1cm of s] (c1) {}; \node[state,right=of c1] (c2) {}; \node[state,right=of c2] (c3) {}; \node[state,accepting,right=of c3] (c4) {};
  \node[state,below right=0.5cm and 1cm of s] (d1) {}; \node[state,accepting,right=of d1] (d2) {};
  \path (s) edge node {$\eps$} (c1) (s) edge node[below left] {$\eps$} (d1)
        (c1) edge node {$a$} (c2) (c2) edge node {$\eps$} (c3) (c3) edge node {$b$} (c4)
        (d1) edge node {$a$} (d2);
  \node[below=2.0cm of l] (l2) {$(ab\cup a)^*$:};
  \node[state,accepting,right=0.3cm of l2] (t) {};
  \node[state,right=0.9cm of t] (s2) {};
  \node[state,above right=0.5cm and 1cm of s2] (e1) {}; \node[state,right=of e1] (e2) {}; \node[state,right=of e2] (e3) {}; \node[state,accepting,right=of e3] (e4) {};
  \node[state,below right=0.5cm and 1cm of s2] (f1) {}; \node[state,accepting,right=of f1] (f2) {};
  \path (t) edge node {$\eps$} (s2) (s2) edge node {$\eps$} (e1) (s2) edge node[below left] {$\eps$} (f1)
        (e1) edge node {$a$} (e2) (e2) edge node {$\eps$} (e3) (e3) edge node {$b$} (e4)
        (f1) edge node {$a$} (f2)
        (e4) edge[bend right=50] node[above] {$\eps$} (s2)
        (f2) edge[bend left=50] node[below] {$\eps$} (s2);
\end{automaton}
در شکل آخر، حالت سمت چپ حالت شروع جدیدِ ساخت ستاره است.
\end{example}

\subsection{از \lr{DFA} به عبارت منظم}

\begin{lemma}\label{lem:dfa-re}
اگر زبانی منظم باشد، یک عبارت منظم آن را توصیف می‌کند.
\end{lemma}

برای این جهت از نوع جدیدی از ماشین استفاده می‌کنیم: \term{ماشین متناهی غیرقطعی تعمیم‌یافته}{Generalized NFA} یا \lr{GNFA}. در \lr{GNFA} برچسب پیکان‌ها می‌تواند \textbf{هر عبارت منظمی} باشد، نه فقط یک نماد یا $\eps$. ماشین با رد شدن از یک پیکان، یک \textbf{بلوک} از نمادهای ورودی را می‌خواند که با برچسب آن پیکان تطبیق دارد.

نقشه اثبات: \lr{DFA} را به یک \lr{GNFA} تبدیل می‌کنیم و بعد حالت‌های \lr{GNFA} را یکی‌یکی حذف می‌کنیم و برچسب‌ها را به‌روز می‌کنیم، تا فقط دو حالت بماند. برچسب تنها پیکان باقی‌مانده همان عبارت منظم مورد نظر است.

\begin{ltrfig}
\begin{tikzpicture}[->,>={Stealth[round]},semithick,every node/.style={draw=jzMain,rounded corners,fill=jzMain!5,font=\small,align=center,minimum height=8mm}]
  \node (a) at (0,0) {3-state\\DFA}; \node (b) at (2.8,0) {5-state\\GNFA}; \node (c) at (5.6,0) {4-state\\GNFA};
  \node (d) at (5.6,-1.6) {3-state\\GNFA}; \node (e) at (2.8,-1.6) {2-state\\GNFA}; \node (f) at (0,-1.6) {regular\\expression};
  \draw (a)--(b); \draw (b)--(c); \draw (c)--(d); \draw (d)--(e); \draw (e)--(f);
\end{tikzpicture}
\end{ltrfig}

برای ساده شدن کار، \lr{GNFA}ها را به یک \textbf{شکل خاص} محدود می‌کنیم:
\begin{itemize}
\item یک حالت شروع و یک حالت پذیرش دارد و این دو با هم فرق دارند.
\item از حالت شروع به هر حالت دیگر پیکان هست، اما هیچ پیکانی وارد حالت شروع نمی‌شود.
\item از هر حالت دیگر به حالت پذیرش پیکان هست، اما هیچ پیکانی از حالت پذیرش خارج نمی‌شود.
\item بین هر دو حالت دیگر، در هر جهت دقیقاً یک پیکان هست و هر حالت یک حلقه روی خودش دارد.
\end{itemize}

\begin{definition}[\lr{GNFA}]
یک \lr{GNFA} پنج‌تایی $(Q,\Sigma,\delta,\qstart,\qacc)$ است که در آن $Q$ مجموعه متناهی حالت‌ها، $\Sigma$ الفبای ورودی، $\qstart$ حالت شروع، $\qacc$ حالت پذیرش و
\[\delta:\big(Q-\{\qacc\}\big)\times\big(Q-\{\qstart\}\big)\to\mathcal{R}\]
تابع گذار است که $\mathcal{R}$ مجموعه همه عبارات منظم روی $\Sigma$ است. یعنی $\delta(q_i,q_j)$ برچسب پیکان از $q_i$ به $q_j$ است.

\lr{GNFA} رشته $w\in\Sigma^*$ را می‌پذیرد اگر بتوان نوشت $w=w_1w_2\cdots w_k$ (هر $w_i\in\Sigma^*$) و دنباله‌ای از حالت‌ها $q_0,q_1,\dots,q_k$ وجود داشته باشد که $q_0=\qstart$، $q_k=\qacc$ و برای هر $i$، $w_i\in L\big(\delta(q_{i-1},q_i)\big)$.
\end{definition}

\textbf{تبدیل \lr{DFA} به \lr{GNFA} با شکل خاص:}
\begin{itemize}
\item یک حالت شروع جدید با پیکان $\eps$ به حالت شروع قدیمی اضافه می‌کنیم.
\item یک حالت پذیرش جدید با پیکان‌های $\eps$ از همه حالت‌های پذیرش قدیمی اضافه می‌کنیم.
\item اگر بین دو حالت چند پیکان (یا یک پیکان با چند برچسب) بود، آن‌ها را با اجتماع یکی می‌کنیم.
\item بین حالت‌هایی که در \lr{DFA} پیکانی نداشتند، پیکانی با برچسب $\emptyset$ می‌گذاریم. (پیکان $\emptyset$ هرگز قابل استفاده نیست، پس زبان تغییر نمی‌کند.)
\end{itemize}

\textbf{حذف یک حالت:} یک حالت $q_{\mathrm{rip}}$ غیر از شروع و پذیرش را انتخاب و حذف می‌کنیم. برای جبران، برچسب هر پیکان $q_i\to q_j$ را طوری عوض می‌کنیم که همه رشته‌هایی را توصیف کند که ماشین را \textbf{یا مستقیماً یا از طریق $q_{\mathrm{rip}}$} از $q_i$ به $q_j$ می‌برند:
\begin{ltrfig}
\begin{tikzpicture}[jzaut,node distance=2cm]
  \node[state] (qi) {$q_i$}; \node[state,right=2.6cm of qi] (qj) {$q_j$};
  \node[state,below right=1.2cm and 0.8cm of qi] (qr) {$q_{\mathrm{rip}}$};
  \path (qi) edge node {$R_4$} (qj) (qi) edge node[below left] {$R_1$} (qr) (qr) edge node[below right] {$R_3$} (qj) (qr) edge[loop below] node {$R_2$} (qr);
  \node at (5.3,-0.5) {\Large$\Longrightarrow$};
  \node[state,right=4.4cm of qj] (pi) {$q_i$}; \node[state,right=3.4cm of pi] (pj) {$q_j$};
  \path (pi) edge node {$(R_1)(R_2)^*(R_3)\cup(R_4)$} (pj);
\end{tikzpicture}
\end{ltrfig}

\begin{algo}{$\textsc{Convert}(G)$}{روی ورودی \lr{GNFA} $G$}
\begin{enumerate}
\item فرض کنید $k$ تعداد حالت‌های $G$ باشد.
\item اگر $k=2$، آنگاه $G$ فقط یک حالت شروع، یک حالت پذیرش و یک پیکان بین آن‌ها با برچسب $R$ دارد. $R$ را برگردان.
\item اگر $k>2$، یک $q_{\mathrm{rip}}\in Q-\{\qstart,\qacc\}$ انتخاب کن و \lr{GNFA} $G'=(Q',\Sigma,\delta',\qstart,\qacc)$ را بساز که $Q'=Q-\{q_{\mathrm{rip}}\}$ و برای هر $q_i\in Q'-\{\qacc\}$ و $q_j\in Q'-\{\qstart\}$:
\[\delta'(q_i,q_j)=(R_1)(R_2)^*(R_3)\cup(R_4),\]
که $R_1=\delta(q_i,q_{\mathrm{rip}})$، $R_2=\delta(q_{\mathrm{rip}},q_{\mathrm{rip}})$، $R_3=\delta(q_{\mathrm{rip}},q_j)$ و $R_4=\delta(q_i,q_j)$.
\item $\textsc{Convert}(G')$ را برگردان.
\end{enumerate}
\end{algo}

\begin{lemma}
برای هر \lr{GNFA} $G$، عبارت $\textsc{Convert}(G)$ با $G$ هم‌ارز است.
\end{lemma}
\begin{proof}
با استقرا روی تعداد حالت‌های $G$، یعنی $k$.

\textbf{پایه} ($k=2$): $G$ فقط یک پیکان از $\qstart$ به $\qacc$ با برچسب $R$ دارد. پس $G$ دقیقاً رشته‌های $L(R)$ را می‌پذیرد و $\textsc{Convert}$ همان $R$ را برمی‌گرداند.

\textbf{گام:} فرض کنید ادعا برای $k-1$ حالت درست باشد و $G$ دارای $k>2$ حالت باشد. کافی است نشان دهیم $G$ و $G'$ هم‌ارزند، چون طبق فرض استقرا $\textsc{Convert}(G')$ با $G'$ هم‌ارز است.

\emph{هر رشته‌ای که $G$ بپذیرد، $G'$ هم می‌پذیرد:} فرض کنید $G$ رشته $w$ را با دنباله حالت‌های $\qstart,q_1,q_2,\dots,\qacc$ بپذیرد. اگر $q_{\mathrm{rip}}$ در این دنباله نباشد، $G'$ هم با همین دنباله $w$ را می‌پذیرد، چون هر برچسب جدید شامل برچسب قدیمی (به عنوان یکی از دو طرف اجتماع، یعنی $R_4$) است. اگر $q_{\mathrm{rip}}$ در دنباله باشد، هر دسته از $q_{\mathrm{rip}}$های پشت سر هم را حذف کنید. اگر این دسته بین $q_i$ و $q_j$ قرار داشته باشد، بخشی از ورودی که از $q_i$ تا $q_j$ خوانده شده با $R_1R_2^*R_3$ تطبیق دارد (یک بار $R_1$، چند بار $R_2$ و یک بار $R_3$). پس $G'$ می‌تواند همین بخش را با پیکان جدید $q_i\to q_j$ بخواند و $w$ را بپذیرد.

\emph{هر رشته‌ای که $G'$ بپذیرد، $G$ هم می‌پذیرد:} هر پیکان $q_i\to q_j$ در $G'$ رشته‌هایی را می‌خواند که $G$ آن‌ها را یا مستقیماً (با $R_4$) یا از طریق $q_{\mathrm{rip}}$ (با $R_1R_2^*R_3$) از $q_i$ به $q_j$ می‌خواند. پس هر محاسبه پذیرنده $G'$ را می‌توان به یک محاسبه پذیرنده $G$ تبدیل کرد.

بنابراین $G$ و $G'$ هم‌ارزند و حکم ثابت است.
\end{proof}

با این لم، اثبات لم \ref{lem:dfa-re} کامل است: برای زبان منظم $A$ یک \lr{DFA} می‌گیریم، آن را به \lr{GNFA} تبدیل می‌کنیم و $\textsc{Convert}$ را اجرا می‌کنیم. و قضیه \ref{thm:re-equiv} از دو لم \ref{lem:re-nfa} و \ref{lem:dfa-re} نتیجه می‌شود.

\begin{example}
\lr{DFA} دوحالته زیر را به عبارت منظم تبدیل می‌کنیم (پیکان‌های $\emptyset$ را نمی‌کشیم):
\begin{automaton}[node distance=1.9cm]
  \node at (-0.9,0) {(a)};
  \node[state,initial] (1) at (0,0) {$1$};
  \node[state,accepting] (2) at (0,-1.8) {$2$};
  \path (1) edge[loop right] node {$a$} (1) (1) edge node[left] {$b$} (2) (2) edge[loop right] node {$a,b$} (2);
  \node at (2.2,0) {(b)};
  \node[state] (s) at (3,0) {$s$};
  \node[state] (1b) at (4.6,0) {$1$};
  \node[state] (2b) at (4.6,-1.8) {$2$};
  \node[state] (ab) at (6.4,-1.8) {$a$};
  \path (s) edge node {$\eps$} (1b) (1b) edge[loop right] node {$a$} (1b) (1b) edge node[left] {$b$} (2b)
        (2b) edge[loop below] node {$a\cup b$} (2b) (2b) edge node {$\eps$} (ab);
  \node at (8.0,0) {(c)};
  \node[state] (sc) at (8.8,0) {$s$};
  \node[state] (1c) at (10.4,0) {$1$};
  \node[state] (ac) at (10.4,-1.8) {$a$};
  \path (sc) edge node {$\eps$} (1c) (1c) edge[loop right] node {$a$} (1c) (1c) edge node[right] {$b(a\cup b)^*$} (ac);
\end{automaton}
در (b) دو حالت جدید $s$ (شروع) و $a$ (پذیرش) اضافه شده‌اند. در (c) حالت $2$ را حذف کرده‌ایم: پیکان جدید $1\to a$ برچسب $b(a\cup b)^*\eps\cup\emptyset=b(a\cup b)^*$ دارد. در آخر با حذف حالت $1$، پیکان $s\to a$ برچسب
\[\eps\,a^*\,b(a\cup b)^*\cup\emptyset=a^*b(a\cup b)^*\]
را می‌گیرد. پس زبان این \lr{DFA} برابر $a^*b(a\cup b)^*$ است: رشته‌هایی که دست‌کم یک $b$ دارند.
\end{example}

\begin{summary}
\begin{itemize}
\item عبارت منظم از $a$، $\eps$ و $\emptyset$ با اجتماع، الحاق و ستاره ساخته می‌شود. اولویت: ستاره، الحاق، اجتماع.
\item $R\cup\emptyset=R$ و $R\eps=R$، اما $R\cup\eps$ و $R\emptyset$ در حالت کلی برابر $R$ نیستند. $\emptyset^*=\{\eps\}$.
\item تبدیل عبارت منظم به \lr{NFA}: ساخت استقرایی با ساخت‌های بستاری.
\item تبدیل \lr{DFA} به عبارت منظم: ساختن \lr{GNFA} و حذف حالت‌ها با فرمول $(R_1)(R_2)^*(R_3)\cup(R_4)$.
\item عبارات منظم، \lr{DFA}ها و \lr{NFA}ها دقیقاً یک دسته زبان را توصیف می‌کنند.
\end{itemize}
\end{summary}
